\Article{217}
For the case where $bb-ac$ is a positive nonsquare number, we showed above that all representations of the number $M$ by the form $app+2bpq+cqq$ (if there are any at all) can be given by one or more formulas such as
\[
p=\tfrac{1}{m}(\mathfrak{A}t+\mathfrak{B}u),\quad q=\tfrac{1}{m}(\mathfrak{C}t+\mathfrak{D}u)
\]
where $\mathfrak{A},\mathfrak{B},\mathfrak{C},\mathfrak{D}$ denote given integers, $m$ the greatest common divisor of the numbers $a,2b,c$, and finally $t,u$ range over all integers satisfying the equation $tt-(bb-ac)uu=mm$. Since all values of $t,u$ can be taken both positively and negatively, we can replace each of those formulas by \textit{four} others:
\[
\begin{aligned}
p&=\tfrac{1}{m}(\mathfrak{A}t+\mathfrak{B}u), &\quad q&=\tfrac{1}{m}(\mathfrak{C}t+\mathfrak{D}u)\\
p&=\tfrac{1}{m}(\mathfrak{A}t-\mathfrak{B}u), &q&=\tfrac{1}{m}(\mathfrak{C}t-\mathfrak{D}u)
\end{aligned}
\]
\SourceLeaf{222}
\[
\begin{aligned}
p&=\tfrac{1}{m}(-\mathfrak{A}t+\mathfrak{B}u),&
q&=\tfrac{1}{m}(-\mathfrak{C}t+\mathfrak{D}u)\\
p&=-\tfrac{1}{m}(\mathfrak{A}t+\mathfrak{B}u),&
q&=-\tfrac{1}{m}(\mathfrak{C}t+\mathfrak{D}u)
\end{aligned}
\]
The total number of formulas is now four times as large as before, while $t$ and $u$ no longer stand for all numbers satisfying the equation $tt-(bb-ac)uu=mm$, but only the positive ones. We must therefore consider each of these forms separately and determine which values of $t$, $u$ give integer values of $x$, $y$.

The formula
\[
p=\tfrac{1}{m}(\mathfrak{A}t+\mathfrak{B}u),\quad
q=\tfrac{1}{m}(\mathfrak{C}t+\mathfrak{D}u)\qquad .\ .\ .\ .\ [1]
\]
gives the following values of $x$, $y$:
\[
x=\frac{\mathfrak{A}t+\mathfrak{B}u+mcd-mbe}{m(bb-ac)},\quad
y=\frac{\mathfrak{C}t+\mathfrak{D}u+mae-mbd}{m(bb-ac)}
\]
But we showed above that all (positive) values of $t$ form a recurrent sequence $t^0,t',t''$ etc., and that the corresponding values of $u$ likewise form a recurrent sequence $u^0,u',u''$ etc.; also, that we can specify a number $\rho$ such that, for any given modulus,
\[
t^\rho\equiv t^0,\quad t^{\rho+1}\equiv t',\quad t^{\rho+2}\equiv t''\ \text{etc.},\quad
u^\rho\equiv u^0,\quad u^{\rho+1}\equiv u'\ \text{etc.}
\]
We take the number $m(bb-ac)$ as this modulus. For brevity, we denote the values of $x$, $y$ obtained by setting $t=t^0$, $u=u^0$, to which we assign the index $0$, by $x^0,y^0$; similarly, those obtained by setting $t=t'$, $u=u'$, by $x',y'$, to which we assign the index $1$, etc. It is then easy to see that if $x^h,y^h$ are integers and $\rho$ is correctly determined, then $x^{h+\rho},y^{h+\rho}$, as well as $x^{h+2\rho},y^{h+2\rho}$ and, in general, $x^{h+k\rho},y^{h+k\rho}$, will also be integers. Conversely, if $x^h$ or $y^h$ is fractional, then $x^{h+k\rho}$ or $y^{h+k\rho}$, respectively, will also be fractional. We readily conclude that if we work out the values of $x$, $y$ with indices $0,1,2\,.\,.\,.\,.\,\rho-1$, and for none of these indices are \textit{both} $x$ \textit{and} $y$ integers, then there is no index at all for which both $x$ and $y$ take integer values. In that case, no integer values of $x$, $y$ can be derived from formula [1]. But if some of those indices, say $\mu,\mu',\mu''$ etc., correspond to integer values of $x$, $y$, then all integer values of $x$, $y$ that can be obtained from formula [1] are those whose indices are contained in one of the formulas $\mu+k\rho,\ \mu'+k\rho,$
\SourceLeaf{223}
$\mu''+k\rho$ etc., where $k$ ranges over all positive integers, including zero.

The remaining formulas containing the values of $p$, $q$ must be treated in exactly the same way. If none of these formulas were to yield integer values of $x$, $y$, the proposed equation could not be solved in integers at all. Whenever it is in fact solvable, however, all its integer solutions can be given by the rules set out above.

\Needspace{8\baselineskip}
\Article{218}
When $bb-ac$ is a square number and $M=0$, all values of $p$, $q$ are contained in two formulas of the kind $p=\mathfrak{A}z$, $q=\mathfrak{B}z$; $p=\mathfrak{A}'z$, $q=\mathfrak{B}'z$, where $z$ ranges over all integers, while $\mathfrak{A},\mathfrak{B},\mathfrak{A}',\mathfrak{B}'$ are given integers, the first and second being relatively prime, and likewise the third and fourth (Article 212). Thus all integer values of $x$, $y$ arising from the first formula are contained in formula [1]
\[
x=\frac{\mathfrak{A}z+cd-be}{bb-ac},\quad
y=\frac{\mathfrak{B}z+ae-bd}{bb-ac}
\]
and all the remaining ones arising from the second formula in this formula [2]:
\[
x=\frac{\mathfrak{A}'z+cd-be}{bb-ac},\quad
y=\frac{\mathfrak{B}'z+ae-bd}{bb-ac}
\]
But since either formula can also give fractional values (unless $bb-ac=1$), we must separate, in each formula, the values of $z$ that make both $x$ and $y$ integers from the others. It is enough, however, to consider the first formula alone, since exactly the same method applies to the other.

Since $\mathfrak{A}$, $\mathfrak{B}$ are relatively prime, we can determine two numbers $\mathfrak{a}$, $\mathfrak{b}$ such that $\mathfrak{a}\mathfrak{A}+\mathfrak{b}\mathfrak{B}=1$. This gives
\[
(\mathfrak{a}x+\mathfrak{b}y)(bb-ac)=z+\mathfrak{a}(cd-be)+\mathfrak{b}(ae-bd)
\]
so it is immediately clear that all values of $z$ capable of producing integer values of $x$, $y$ must be congruent to the number $\mathfrak{a}(be-cd)+\mathfrak{b}(bd-ae)$ modulo $bb-ac$, or be contained in the formula $(bb-ac)z'+\mathfrak{a}(be-cd)+\mathfrak{b}(bd-ae)$, where $z'$ ranges over all integers. From this we easily obtain, in place of formula [1], the following:
\SourceLeaf{224}
\[
\begin{aligned}
x&=\mathfrak{A}z'+\mathfrak{b}\times\frac{\mathfrak{A}(bd-ae)-\mathfrak{B}(be-cd)}{bb-ac}\\
y&=\mathfrak{B}z'-\mathfrak{a}\times\frac{\mathfrak{A}(bd-ae)-\mathfrak{B}(be-cd)}{bb-ac}
\end{aligned}
\]
It clearly gives, either for all values of $z'$ or for none, integer values of $x,y$. The first case always occurs when $\mathfrak{A}(bd-ae)$ and $\mathfrak{B}(be-cd)$ are congruent modulo $bb-ac$, and the second when they are incongruent. --- Formula [2] must be treated in exactly the same way, separating its integer solutions (if it can give any) from the others.

\Needspace{8\baselineskip}
\Article{219}
When $bb-ac=0$, the form $axx+2bxy+cyy$ can be written as $m(\alpha x+\beta y)^2$, where $m,\alpha,\beta$ are integers (Article 215). Set $\alpha x+\beta y=z$; the proposed equation becomes
\[
mzz+2dx+2ey+f=0
\]
and from this and $z=\alpha x+\beta y$ we obtain
\[
x=\frac{\beta mzz+2ez+\beta f}{2\alpha e-2\beta d},\quad
y=\frac{\alpha mzz+2dz+\alpha f}{2\beta d-2\alpha e}
\]
It is now clear that, unless $\alpha e=\beta d$ (a case we will consider separately in a moment), the values of $x$, $y$ derived from these formulas by giving $z$ any value satisfy the proposed equation. It remains only to show how to determine the values of $z$ that give integer values of $x$, $y$.

Since $\alpha x+\beta y=z$, only \textit{integer} values can be taken for $z$. Moreover, if some value of $z$ makes both $x$ and $y$ integers, all values of $z$ congruent to it modulo $2\alpha e-2\beta d$ clearly also produce integer values. So substitute for $z$ all integers from $0$ to $2\alpha e-2\beta d-1$ (when $\alpha e-\beta d$ is positive) or to $2\beta d-2\alpha e-1$ (when $\alpha e-\beta d$ is negative), inclusive. If none of these values makes both $x$ and $y$ integers, then no value of $z$ at all will produce integer values of $x,y$, and the proposed equation cannot be solved in integers. But if some of those values of $z$ give integer values of $x,y$, say $\zeta,\zeta',\zeta''$ etc. (which can also be found by solving quadratic congruences using the principles of Section IV), \textit{all} solutions follow\SourceLeaf{225} by setting $z=(2\alpha e-2\beta d)v+\zeta$, $z=(2\alpha e-2\beta d)v+\zeta'$ etc., where $v$ ranges over all integers.

\Needspace{8\baselineskip}
\Article{220}
For the excluded case where $\alpha e=\beta d$, we must find a separate method. Suppose $\alpha$, $\beta$ are relatively prime, as Article 215, I shows we may assume. Then $\frac{d}{\alpha}=\frac{e}{\beta}$ is an integer (Article 19), which we set $=h$. The proposed equation now takes the form
\[
(m\alpha x+m\beta y+h)^2-hh+mf=0
\]
and so clearly has no rational solution unless $hh-mf$ is a square number. Let $hh-mf=kk$; the following two equations are clearly equivalent to the proposed equation:
\[
m\alpha x+m\beta y+h+k=0,\quad\text{and}\quad
m\alpha x+m\beta y+h-k=0
\]
\textit{i.e.}, every solution of the proposed equation also satisfies one or the other of these equations, and conversely. The first equation clearly cannot be solved in integers unless $h+k$ is divisible by $m$, and similarly the second admits no integer solution unless $h-k$ is divisible by $m$. These conditions also suffice for the solvability of each equation (since $\alpha$, $\beta$ are assumed relatively prime), and all the solutions can be given by the known rules.

\Needspace{8\baselineskip}
\Article{221}
We illustrate the case considered in Article 217 (since it is the most difficult of all) by an example. Let the proposed equation be $xx+8xy+yy+2x-4y+1=0$. From this, by first introducing the other unknowns $p=15x-9$. $q=15y+6$, we derive the equation $pp+8pq+qq=-540$. All its integer solutions are found to be contained in the following four formulas:
\[
\begin{aligned}
p&=6t,&q&=-24t-90u\\
p&=6t,&q&=-24t+90u\\
p&=-6t,&q&=24t-90u\\
p&=-6t,&q&=24t+90u
\end{aligned}
\]
where $t$, $u$ range over all positive integers satisfying the equation
\SourceLeaf{226}
$tt-15uu=1$, which are contained in the formula
\[
\begin{aligned}
t&=\tfrac{1}{2}\bigl((4+\sqrt{15})^n+(4-\sqrt{15})^n\bigr)\\
u&=\frac{1}{2\sqrt{15}}\bigl((4+\sqrt{15})^n-(4-\sqrt{15})^n\bigr)
\end{aligned}
\]
when $n$ ranges over all positive integers (including zero). Thus all values of $x,y$ are contained in the following formulas:
\[
\begin{aligned}
x&=\tfrac{1}{5}(2t+3),&y&=-\tfrac{1}{5}(8t+30u+2)\\
x&=\tfrac{1}{5}(2t+3),&y&=-\tfrac{1}{5}(8t-30u+2)\\
x&=\tfrac{1}{5}(-2t+3),&y&=\tfrac{1}{5}(8t-30u-2)\\
x&=\tfrac{1}{5}(-2t+3),&y&=\tfrac{1}{5}(8t+30u-2)
\end{aligned}
\]
Applying our rules correctly, we find that to obtain \textit{integer} values, the first and second formulas require those values of $t,u$ that arise from an \textit{even} index $n$; the third and fourth require those obtained from an \textit{odd} $n$. --- The simplest solutions are $x=1,-1,-1$; $y=-2,0,12$, respectively.

We should also note that the solution of the problem explained in the preceding articles can usually be shortened by various devices, especially in excluding unusable solutions, \textit{i.e.}, those involving fractions. But to avoid going on too long, we have had to pass over these here.
